\documentclass[11pt,a4paper]{article}
\usepackage[T1]{fontenc}
\usepackage[utf8]{inputenc}
\usepackage[main=italian,provide=*]{babel}
\usepackage{amsmath,amssymb}
\usepackage{geometry}
\geometry{margin=2.3cm,top=2.5cm,bottom=2.7cm}
\usepackage{booktabs}
\usepackage[table]{xcolor}
\usepackage{tcolorbox}
\tcbuselibrary{skins,breakable}
\usepackage{titlesec}
\usepackage{enumitem}
\usepackage{needspace}
\usepackage{fancyhdr}
\usepackage{hyperref}
\hypersetup{colorlinks=true,linkcolor=Sepia,citecolor=Sepia,urlcolor=Sepia}
\newcommand{\Tr}{\operatorname{Tr}}
\newcommand{\ket}[1]{\lvert #1 \rangle}
\newcommand{\bra}[1]{\langle #1 \rvert}

% --- palette ---
\definecolor{inkblue}{HTML}{1B3A5C}
\definecolor{lightblue}{HTML}{EAF1F8}
\definecolor{accent}{HTML}{B0562A}
\definecolor{softgray}{HTML}{6B6B6B}
\definecolor{okgreen}{HTML}{2E6B3E}
\definecolor{okgreenbg}{HTML}{EAF5EC}

% --- titoli di sezione ---
\titleformat{\section}
  {\normalfont\Large\bfseries\color{inkblue}}
  {\thesection}{0.6em}{}[{\color{inkblue!40}\titlerule}]
\titlespacing{\section}{0pt}{0.7em}{0.4em}
\titleformat{\subsection}
  {\normalfont\large\bfseries\color{accent}}
  {}{0em}{}
\titlespacing{\subsection}{0pt}{0.5em}{0.2em}

\pagestyle{fancy}
\fancyhf{}
\renewcommand{\headrulewidth}{0.4pt}
\fancyhead[L]{\small\color{softgray}VQE noise-aware --- estensione del VQE con termine DM sotto rumore}
\fancyhead[R]{\small\color{softgray}\thepage}
\fancyfoot[C]{\small\color{softgray}Tesi triennale, Samuele --- Relatore: Prof.\ Alessandro Chiesa}

% --- box riutilizzabili ---
\newtcolorbox{keyeq}[1][]{
  colback=lightblue, colframe=inkblue!70, boxrule=0.6pt,
  arc=2pt, left=8pt, right=8pt, top=3pt, bottom=3pt, #1
}
\newtcolorbox{okbox}[1]{
  colback=okgreenbg, colframe=okgreen!60, boxrule=0.6pt,
  arc=2pt, left=7pt, right=7pt, top=2pt, bottom=2pt,
  fonttitle=\bfseries\color{okgreen}, title={#1}
}

\title{\vspace{-1.5em}\color{inkblue}VQE noise-aware --- estensione del VQE con termine DM sotto rumore\\[0.2em]
\Large\normalfont\color{softgray}Rioptimizzazione sotto rumore e controllo su $N^*$ della quantum simulation (Trotter) sotto rumore}
\author{Note di lavoro --- Tesi triennale, Samuele\\ Relatore: Prof.\ Alessandro Chiesa}
\date{}

\begin{document}
\sloppy
\maketitle
\thispagestyle{fancy}
\vspace{-2.2em}

File: \texttt{vqe\_noise\_aware\_dimero.py}, \texttt{validate\_vqe\_noise\_aware\_dimero.py}.
Punto di lavoro: $b/J=0.35$, $D/J=0.80$, $J=1$ (``test 2''). Modello di
rumore: \texttt{noise\_model\_dimero.py} (modello di rumore, calibrazione
\texttt{ibm\_torino}).

\section*{Scope dichiarato}
\addcontentsline{toc}{section}{Scope dichiarato}

\begin{itemize}[leftmargin=1.3em,itemsep=0.3em]
\item Questa è un'estensione \emph{facoltativa} del VQE con termine DM sotto rumore,
non uno step mancante della pipeline richiesta (già chiusa in tutti e
cinque i passi, vedi \texttt{log\_decisioni.md}).
\item Si rioptimizza \textbf{solo} il VQE con termine DM sotto rumore.
\item Un \textbf{solo} controllo a valle: si riverifica se questo sposta
$N^*=8$ della quantum simulation (Trotter) sotto rumore
(\texttt{trotter\_rumoroso\_dimero.py}), riusando i
suoi stessi ingredienti con i nuovi parametri al posto di quelli ideali.
\item I correlatori dinamici sotto rumore e lo scan sui parametri di
rumore \textbf{non} vengono ritoccati qui: l'estensione a questi due
stadi si basa su un argomento di struttura, non su una
riverifica numerica diretta (vedi sotto, punto (b)--(d)) --- criterio di
proporzionalità dato il tempo disponibile prima della sessione di laurea
di settembre 2026.
\end{itemize}

\section*{Metodologia}
\addcontentsline{toc}{section}{Metodologia}

\begin{itemize}[leftmargin=1.3em,itemsep=0.35em]
\item \textbf{Due domande diverse (richiamo dalla versione originale).}
Il VQE con termine DM sotto rumore misura la \emph{degradazione} di un risultato ideale: $\theta^*$ è
ottimizzato \textbf{senza} rumore (Parte 1), poi valutato \textbf{con}
rumore. Qui si fa la domanda del VQE \emph{noise-aware}: rifare l'intero
ciclo di ottimizzazione con il rumore acceso fin dall'inizio, come in
un'esecuzione reale su hardware dove non esiste una versione ideale a cui
accedere per calibrare l'ottimizzazione.

\item \textbf{(a) Argomento fisico per la previsione.} Il canale
depolarizzante è isotropo: per un canale così, l'azione su uno stato puro
si scrive $\rho(\theta)=(1-\lambda)\ket{\psi(\theta)}\bra{\psi(\theta)}+\lambda\,\mathbf{1}/d$
(schematicamente), quindi
\begin{keyeq}
\centering
$E_\text{rumoroso}(\theta)=\Tr[\rho(\theta)H] = (1-\lambda)\,E_\text{ideale}(\theta) + \lambda\,\Tr[H]/d$
\end{keyeq}
è una trasformazione \emph{affine} di $E_\text{ideale}(\theta)$ in
$\theta$: stesso segno di $\lambda>0$, quindi lo stesso $\theta^*$ che
minimizza $E_\text{ideale}$ minimizza anche $E_\text{rumoroso}$. Inoltre
il conteggio CNOT del circuito transpilato non dipende dai valori
numerici dei parametri (verificato sotto): nessuna direzione del
paesaggio verso cui convenga spostarsi per ``risparmiare'' rumore.
\textbf{Previsione}: $\theta^*_\text{rumore}\approx\theta^*_\text{ideale}$,
nessun vantaggio misurabile a rioptimizzare.

\item \textbf{(b) Perché questo si estende ai correlatori dinamici e allo
scan sui parametri senza riverificarli lì.} Già dimostrato nella
quantum simulation (Trotter) sotto rumore che il contributo della
preparazione è una costante additiva indipendente da $N$: non entra né
nel termine che decresce con $N$ (Trotter) né in quello che cresce con $N$
(rumore accumulato). I correlatori dinamici e lo scan sui parametri
condividono la stessa struttura con la
stessa preparazione, quindi lo stesso argomento si applica identico a
tutti e tre.

\item \textbf{(c) Perché il controllo numerico resta comunque
necessario.} (a) è una previsione, non un fatto acquisito: se il
controllo su $N^*=8$ (quantum simulation Trotter sotto rumore) la
conferma, l'estensione ai correlatori dinamici e allo scan sui parametri è
giustificata dallo stesso argomento; se la smentisse, si riaprirebbero
anche quelli, perché la premessa che li chiudeva sarebbe caduta.

\item \textbf{(d) Criterio di proporzionalità.} Ripetere lo stesso
controllo tre volte quando il meccanismo è identico non aggiunge
comprensione, solo costo --- rilevante data la scadenza di settembre.

\item \textbf{Nota tecnica --- una trappola scoperta e corretta.} Un primo
tentativo transpilava l'ansatz \textbf{una sola volta in forma parametrica}
(parametri simbolici), riusando poi \texttt{assign\_parameters} ad ogni
valutazione --- lo stesso principio che funziona per il blocco di Trotter
ripetuto $N$ volte (quantum simulation Trotter sotto rumore).
Verificato che qui è \textbf{sbagliato}: con
parametri simbolici il transpilatore non può fondere le rotazioni a 1
qubit che dipendono dal valore numerico di $\theta$, producendo un
circuito molto più rumoroso (19 \texttt{rz} + 14 \texttt{sx} contro 11
\texttt{rz} + 8 \texttt{sx}, a parità di 2 CNOT). La scelta corretta,
coerente con \texttt{vqe\_dm\_rumoroso\_dimero.py}, è transpilare
\textbf{dopo} l'assegnazione dei parametri, ad ogni valutazione
dell'obiettivo: più costoso per valutazione ma l'unico che riproduce il
conteggio di gate minimo già stabilito nella pipeline.
\end{itemize}

\section*{Verifica}
\addcontentsline{toc}{section}{Verifica}

\begin{okbox}{Cinque controlli superati}
\begin{itemize}[leftmargin=1.2em,itemsep=0.2em]
\item \textbf{Invarianza CNOT}: CNOT$=2$ su 5 assegnazioni casuali di
$\theta$ (11 \texttt{rz} + 8 \texttt{sx} + 2 \texttt{cx} in ciascuna) ---
premessa fisica di (a) confermata.
\item \textbf{Limite di rumore nullo}: VQE noise-aware con rumore spento
ritrova $E_0=-3.57321145$ (registrato, Parte 1) entro $1.3\times10^{-12}$.
\item \textbf{Robustezza del risultato rumoroso}: un multistart
indipendente più ampio ($R=20$, nessun seme) converge allo stesso ottimo
rumoroso trovato con $R=6$+seme entro $2.6\times10^{-6}$ --- non è un
artefatto del seme. Nota onesta: con solo $R=6$ e nessun seme (stesso $R$
della versione originale) il minimo trovato può restare a distanza
$10^{-4}$--$10^{-2}$ dall'ottimo, a seconda del seed casuale (dispersione
reale del paesaggio energetico, caratterizzata su 8 seed) --- per questo
il risultato principale include il seme dei parametri ideali fra i punti
di partenza.
\item \textbf{Confronto con la versione originale (VQE+DM rumoroso)}:
$|\Delta E|=1.16\times10^{-7}$,
$|\Delta F|=1.76\times10^{-8}$ rispetto al riuso dei parametri ideali sotto
rumore --- entrambi ben sotto la soglia $10^{-4}$ attesa.
\item \textbf{$N^*$ della quantum simulation Trotter sotto rumore}:
invariato, $N^*=8$ con i parametri
noise-aware (era $N^*=8$ con i parametri ideali).
\end{itemize}
\end{okbox}

\section*{Risultato: confronto energia e fedeltà}
\addcontentsline{toc}{section}{Risultato}

\begin{center}
\renewcommand{\arraystretch}{1.25}
\begin{tabular}{@{}lcc@{}}
\toprule
\textbf{Quantità} & \textbf{Riuso $\theta^*_\text{ideale}$ (VQE+DM rumoroso)} & \textbf{VQE noise-aware}\\
\midrule
$E$ (rumoroso) & $-3.50164605$ & $-3.50164616$\\
$F$ (rumorosa) & $0.98500601$ & $0.98500602$\\
CNOT & $2$ & $2$\\
\bottomrule
\end{tabular}
\end{center}

\begin{center}
\renewcommand{\arraystretch}{1.2}
\begin{tabular}{@{}cc@{}}
\toprule
\textbf{Quantità} & \textbf{Valore}\\
\midrule
$\Delta E$ (noise-aware $-$ riuso ideale) & $-1.16\times10^{-7}$\\
$\Delta F$ (noise-aware $-$ riuso ideale) & $+1.76\times10^{-8}$\\
\bottomrule
\end{tabular}
\end{center}

\noindent Scostamento centinaia di migliaia di volte più piccolo della
scala fisica del problema (perdita di fedeltà $\sim1.5\times10^{-2}$ a
rumore acceso): rioptimizzare sotto rumore non porta alcun vantaggio
misurabile al punto di lavoro test 2. Conferma diretta dell'argomento (a):
il canale isotropo rende il paesaggio energetico rumoroso una
trasformazione affine di quello ideale, con lo stesso minimizzatore.

\section*{Risultato: controllo su $N^*$ (quantum simulation Trotter sotto rumore)}
\addcontentsline{toc}{section}{Controllo su $N^*$}

\begin{center}
\renewcommand{\arraystretch}{1.2}
\begin{tabular}{@{}lcc@{}}
\toprule
\textbf{Parametri usati come preparazione} & $\boldsymbol{N^*}$ & $\boldsymbol{F(N^*)}$\\
\midrule
$\theta^*_\text{ideale}$ (versione originale, Trotter rumoroso) & $8$ & $0.817725$\\
$\theta^*_\text{noise-aware}$ (qui) & $8$ & $0.817767$\\
\bottomrule
\end{tabular}
\end{center}

\begin{okbox}{$N^*$ non si sposta}
Il massimo di $F(N)$ resta a $N^*=8$, con lo stesso valore di fedeltà
entro $4\times10^{-5}$. La previsione (b)--(c) è quindi confermata
numericamente, non solo argomentata: il contributo di preparazione resta
il termine costante (floor) già caratterizzato nella quantum simulation
Trotter sotto rumore, indipendente
da \emph{quale} versione dei parametri VQE (ideali o noise-aware) lo
genera.
\end{okbox}

\needspace{8\baselineskip}
\section*{Conclusione ed estensione ai correlatori dinamici e allo scan sui parametri}
\addcontentsline{toc}{section}{Conclusione}

Il controllo numerico richiesto (punto (c) della metodologia) conferma la
previsione (a): per un canale depolarizzante isotropo, rioptimizzare
l'ansatz PMA-2q$\cdot$3 sotto rumore non cambia in modo misurabile né
l'energia/fedeltà del VQE+DM rumoroso né la posizione di $N^*$ della
quantum simulation Trotter sotto rumore. Per l'argomento di struttura
(b), la stessa conclusione si estende ai correlatori dinamici sotto
rumore ($N^*=5$) e allo scan sui parametri di rumore, che
condividono l'identica preparazione VQE e l'identica struttura additiva
$\varepsilon_\text{tot}(N)\simeq a/N+g\varepsilon_{2q}N+\delta_\text{prep}$
--- non riverificati numericamente qui per il criterio di
proporzionalità (d), ma la premessa che li giustificava (il floor di
preparazione è costante e indipendente da $N$) non è stata falsificata da
questo controllo, anzi è stata rinforzata.

\end{document}
